\section{Related Work}

\subsection{Transformer-to-SSM Distillation}

Converting quadratic Transformers to linear-time State Space Models has emerged as a practical path to efficient long-context processing. MOHAWK \citep{bick2024mohawk} introduces a three-stage progressive distillation approach for Phi-1.5 to Phi-Mamba conversion, achieving 85\% of teacher performance with 3B tokens. Stage 1 aligns attention mixer outputs directly via $\|$TeacherMixer - StudentMixer$\|$ minimization---a matrix-level objective that preserves attention patterns.

CAB \citep{wang2025cab} takes a different approach, aligning token-level representations. Query and key projections (Q, K) map to SSM parameters (B, C) through learned MLP bridges, avoiding O($L^2$) attention map materialization. This architectural efficiency enables longer sequence training but raises questions about what information is transferred.

MambaInLlama \citep{wang2024mammainllama} explores hybrid conversion, replacing attention layers with Mamba blocks progressively. T2MD \citep{zhang2024t2md} extends distillation to multimodal settings. Both focus on architectural choices rather than comparing supervision signals.

\textbf{Gap}: No prior work systematically compares matrix-level versus token-level distillation objectives under controlled conditions, particularly across sequence lengths.

\subsection{Length Generalization in Transformers}

Position encoding limits Transformer length generalization. Learned absolute positions (as in Phi-1.5, GPT-2) cannot extrapolate beyond training length---the position embedding matrix has fixed size. Rotary Position Embeddings (RoPE) \citep{su2021roformer} and ALiBi \citep{press2021alibi} enable some extrapolation by encoding relative positions.

Position interpolation \citep{chen2023extending} and NTK-aware scaling extend RoPE-based models to longer contexts. These techniques operate on attention mechanisms directly, not on distillation objectives.

\textbf{Gap}: Length generalization research focuses on attention computation, not on how distillation objective choice affects generalization to the student architecture.

\subsection{Hybrid Architectures}

Recent work explores mixing attention and linear layers. Jamba \citep{lieber2024jamba} interleaves Mamba and attention layers. RWKV \citep{peng2023rwkv} achieves linear complexity through gated RNN-like structures. Liger \citep{lan2025liger} repurposes key matrix weights for gating, achieving 93\% performance recovery with zero new parameters.

\textbf{Gap}: Hybrid architecture research optimizes layer composition. We address a complementary question: given a target architecture, which distillation objective produces representations that generalize across lengths?

\subsection{Knowledge Distillation}

Knowledge distillation \citep{hinton2015distilling} transfers knowledge from larger to smaller models via soft label matching. Feature-based distillation \citep{romero2015fitnets} aligns intermediate representations. Attention transfer \citep{zagoruyko2017attention} specifically matches attention maps.

For architecture conversion (Transformer to SSM), the ``teacher'' and ``student'' have fundamentally different inductive biases. Matrix-level objectives (like MOHAWK Stage 1) resemble attention transfer; token-level objectives (like CAB) resemble feature distillation on projections rather than aggregated maps.

Our work positions distillation objective selection as a design choice with measurable consequences for length generalization---a dimension not explored in standard distillation literature.
